\scanpage{118}
\uline{Dimension modulaire 1} Tous les sommets de $G$, sauf un seul $s_0$, sont de genre $0$ et d'ordre $3$, tandis que $s_0$ est soit de genre $0$, ordre $4$, soit de genre $1$, ordre $1$.

\figure{\begin{tikzpicture}[scale=.8]
\fill (0,0) circle (2pt);
\draw (0,0)--(-.8,.38) (0,0)--(-.4,.72) (0,0)--(.37,.72) (0,0)--(.95,.63);
\draw[densely dotted] (-1.55,.48)--(-.92,.55) (-1.05,.92)--(-.5,.96) (-.21,.96)--(.36,.92) (.82,.85)--(1.33,.91);
\node[below] at (0,0) {$s_0$};
\draw[densely dotted] (3.8,.3)--(4.4,.3);
\draw (4.7,.32)--(6.2,.4);
\fill (6.2,.4) circle (2pt);
\node[below] at (6.2,.4) {$1$};\node[right] at (6.45,.25) {$s_0$};
\end{tikzpicture}}

Traitons d'abord le premier cas. Soit $I=\widehat{\vec A}_{s_0}$, ens. de cardinal $4$. On a un iso canonique
\[
M_G\simeq M_{0,I}\qquad
(\simeq\overbrace{\mathbb{P}^1_{\mathbb{C}}\setminus\{0,1,\infty\}}^{\mathbb{U}_{0,3}},
\]
\margin{isomorphisme qui dépend d'une numérotation des éléments de $I$}
\[
\deleted{\unsure{SP_G}\simeq}
\]
La donnée d'une courbe $X$ de type $G$ équivaut à celle d'une \deleted{\unsure{sphère}} droite projective \inserted{$Y$} $I$-marquée.

On a
\[
\begin{aligned}
M_{0,I}&=\deleted{M_{0,4}\wedge}\operatorname{Rep}(I)\wedge_{\unsure{\mathfrak{S}_4|\mathfrak{G}_4}}M^!_{0,4}\\
&=\operatorname{Rep}(I)\wedge_{\unsure{\mathfrak{S}_4|\mathfrak{G}_4}}\mathbb{U}_{0,3}
\end{aligned}
\]
\margin{NB $\operatorname{Rep}(I)=\operatorname{Bij}(I_4,I)$ avec $I_4=\mathbb{N}\cap[1,4]$}
où $\unsure{\mathfrak{S}_4|\mathfrak{G}_4}$ opère sur $\mathbb{U}_{0,3}$ par l'intermédiaire de l'action du groupe $\unsure{\mathfrak{S}_3|\mathfrak{G}_3}$ et de l'homomorphisme
\diagram{\begin{tikzcd}
\unsure{\mathfrak{S}_4|\mathfrak{G}_4}\arrow[r] & \unsure{\mathfrak{S}_3|\mathfrak{G}_3}
\end{tikzcd}}
où $\unsure{\mathfrak{S}_3|\mathfrak{G}_3}$ est identifié au \deleted{\unsure{l'ensemble}} \inserted{groupe} des permutations de l'ens des partitions de type $2,2$ de l'ens à $4$ éléments $[1,4]\cap\mathbb{N}=I_4$, ces partitions étant numérotées de $1$ à $3$, la partition \uline{$P_i$} ($1\leq i\leq3$) étant \inserted{formée de} $\{i,4\}$ et de son complémentaire \deleted{$\unsure{I_4}$} $I_3\setminus\{i\}$ dans $I_4$. On aura, si $T_i$ ($i\in I$) sont les fibrés tangents unitaires de $Y$ aux pts de $I=\widehat{\vec A}_{s_0}$, pour \deleted{$\unsure{P}$} $a\in A$ :
\mspage{117}
